%% Master-CV - Claudio Swijsen
%% Volledig overzicht van competenties, ervaring en realisaties.
%% Gebruik als bronbestand bij het bouwen van rolspecifieke CV's (main_<bedrijf>_<rol>.tex).
%% Compileren met: lualatex main_example.tex
%% (pdflatex faalt vaak op moderne MiKTeX met fontawesome5 font-expansion errors.)

\documentclass[11pt,a4paper,sans]{moderncv}
\moderncvstyle{banking}
\moderncvcolor{blue}

% Force both first and last name AND section headings to render in moderncv
% blue (color1). Default banking on lualatex+MiKTeX leaves these black, which
% looks inconsistent with the rest of the blue accent scheme.
\renewcommand*{\firstnamestyle}[1]{{\fontsize{34}{36}\bfseries\upshape\color{color1}#1}}
\renewcommand*{\lastnamestyle}[1]{{\fontsize{34}{36}\bfseries\upshape\color{color1}#1}}
\renewcommand*{\sectionstyle}[1]{{\sectionfont\color{color1}#1}}

\usepackage[utf8]{inputenc}
\usepackage{hyperref}
\hypersetup{
    colorlinks=true,
    linkcolor=blue,
    filecolor=magenta,
    urlcolor=blue,
    pdftitle={Claudio Swijsen - CV},
    pdfpagemode=FullScreen,
}
\usepackage[scale=0.80]{geometry}
\usepackage{import}
\usepackage{needspace}

% personal data
\name{Claudio}{Swijsen}
\address{Lummen, Belgi\"e}{}{}
\phone[mobile]{+32~491~62~16~21}
\email{cswijsen@gmail.com}
\extrainfo{\href{https://linkedin.com/in/claudioswijsen}{LinkedIn} (linkedin.com/in/claudioswijsen), \href{https://andthenwhat.be}{andthenwhat.be}}

\begin{document}

\makecvtitle

% ============================================================
%     PROFIEL
% ============================================================

\vspace{6pt}
\small{Wie de klantnood niet scherp heeft, optimaliseert de verkeerde dingen. Klantnoden scherp krijgen en vertalen naar proposities: dat is mijn vak. Twintig jaar B2B over media, automotive, tech en energie, telkens op de plek waar de propositie het probleem was. Outside-in, met behavioural design als vast instrument. Praktijk op het niveau dat een masterdiploma certificeert, aangescherpt aan INSEAD en Vlerick.}

\vspace{4pt}
\small{\textbf{In cijfers:} groeistrategie tot 2030 goedgekeurd door de raad van bestuur \textperiodcentered\ fleet sales +56\% in \'e\'en jaar \textperiodcentered\ HIP, voorloper van het cafetariaplan, uitgegroeid tot 500+ wagens/jaar \textperiodcentered\ B2B-medium met 110.000 bezoekers/maand}

% ============================================================
%     KERNCOMPETENTIES
% ============================================================

\section{Kerncompetenties}
\vspace{1pt}
\begin{itemize}

\item \textbf{Proposities en strategie}: value proposition-ontwikkeling, waardeproposities (Playing to Win, Blue Ocean), messaging, storytelling, go-to-market, businessmodelstrategie.

\item \textbf{Klantinzicht en data}: persona's en kwalitatief klantonderzoek, marktdata-analyse en datavisualisatie, behavioural design.

\item \textbf{Prototyping}: bouwt zelf werkende prototypes in Next.js, React en HTML om idee\"en snel tastbaar en testbaar te maken.

\item \textbf{AI}: AI-native marketingworkflows met Claude Code, Leading AI \& digital marketing strategy (INSEAD), zes Anthropic-certificaten (2026).

\item \textbf{Leiderschap}: leidde marketing- en salesteams tot 11 medewerkers; C-level sparringpartner van CEO tot raad van bestuur.

\end{itemize}

% ============================================================
%     PROFESSIONELE ERVARING
% ============================================================

\section{Professionele ervaring}
\vspace{3pt}
\begin{itemize}

% --- Comfort Energy Group ---
\item{\cventry{2024--2026}{Strategic Growth Leader}{Comfort Energy Group}{Belgi\"e}{}{\vspace{1pt}
\begin{itemize}
    \item Schreef de winning sales propositions, de commerci\"ele invulling van de how to win: proposities die klantnoden invullen die de concurrentie laat liggen; de uitrol loopt.
    \item Herkende in grootschalig klantonderzoek en eigen marktdata het patroon achter structureel dalende volumes en vertaalde dat naar het inzicht dat de hele strategie draagt: klantlevenswaarde boven transactie.
    \item Ontwikkelde die groeistrategie samen met de CEO volgens het Playing to Win-framework, verdeeld over vier fiscale jaren tot 2030; goedgekeurd door de raad van bestuur.
    \item End-to-end verantwoordelijk voor de strategic growth marketing; bouwde de proposities over alle afdelingen heen, samen met CDO, CFO, Chief M\&A Officer, Chief Supply Chain Officer en Operations.
    \item Bracht de Belgische markt met eigen data-analyse tot op straatniveau in kaart; die kaart stuurt de where to play en de proposities.
    \item Aangeduid als creatieve lead van het nieuwe transformation office.
\end{itemize}}}

\vspace{3pt}

% --- Do Don't Try ---
\item{\cventry{2023--2024}{Co-founder en Managing Director}{Do Don't Try, strategische agency}{Hasselt}{}{\vspace{1pt}
\begin{itemize}
    \item Richtte een all-round agency op, gebouwd op behavioural design; leverde propositietrajecten voor een internationale B2B-speler en voor startende en gevestigde ondernemers, van persona-onderzoek tot go-to-market.
\end{itemize}}}

\vspace{3pt}

% --- (G) Inspirational ---
\item{\cventry{2022--2024}{Managing Director}{(G) Inspirational, vennootschap boven het ginmerk Holy Water}{Diest}{}{\vspace{1pt}
\begin{itemize}
    \item Nam de operationele leiding van een startup in moeilijkheden (meer dan 350.000 euro verlies): sneed in de kostenbasis, herzag prijs en distributie, zette een B2C-webshop op en besliste eind 2024 te stoppen toen stoppen de juiste keuze was.
\end{itemize}}}

\vspace{3pt}

% --- A&M Group Head of Marketing ---
\needspace{5\baselineskip}
\item{\cventry{2022}{Head of Marketing}{A\&M Group, automotive groep}{Hasselt}{}{\vspace{1pt}
\begin{itemize}
    \item Gepromoveerd tot head of marketing voor de volledige groep; leverde de rebranding die alle dealergroepen onder \'e\'en identiteit bracht, van strategie en naamgeving tot uitrol.
\end{itemize}}}

\vspace{3pt}

% --- A&M / Groep Delorge ---
\item{\cventry{2020--2022}{B2B- en Marketingmanager}{A\&M Group / Groep Delorge, automotive groep}{Hasselt}{}{\vspace{1pt}
\begin{itemize}
    \item Twee voltijdse functies tegelijk: leidde twee teams, 4 marketeers en 7 medewerkers in fleet, voor 53 showrooms en 8 merken (VW-groep, BMW, MINI, JLR).
    \item Groeide de fleet sales van 2.850 wagens (2020) naar 4.450 (2021), een stijging van 56\% op \'e\'en jaar.
\end{itemize}}}

\vspace{3pt}

% --- i3-Technologies ---
\item{\cventry{2018--2020}{Director of Sales}{i3-Technologies, onderwijs- en corporate technologie}{internationaal}{}{\vspace{1pt}
\begin{itemize}
    \item Leidde 4 regional sales managers (Benelux, Frankrijk, Nordics, Baltics, DACH); groeide de omzet in 2019 van 13 naar 15 miljoen euro (+15\%).
\end{itemize}}}

\vspace{3pt}

% --- One Consultants ---
\item{\cventry{2017--2018}{Marketing- en Business Development Manager}{One Consultants, SAP HANA-consultancy}{Europa}{}{\vspace{1pt}
\begin{itemize}
    \item Opende de deur bij SAP Belgi\"e; dat partnership mondde uit in een spreekplek op een Europees SAP-klantenevent in Parijs, een podium dat anders enkel SAP zelf betreedt.
    \item Leverde de rebranding.
\end{itemize}}}

\vspace{3pt}

% --- Groep Delorge B2B ---
\item{\cventry{2015--2017}{B2B-Manager}{Groep Delorge}{Hasselt}{}{\vspace{1pt}
\begin{itemize}
    \item Verhoogde de fleet sales met 20\% in 2016.
    \item Lanceerde de innovatie HIP (Happiness as Incentive Plan), voorloper van het cafetariaplan: 60 wagens in de markt, na 2017 doorgegroeid tot 500+ per jaar.
\end{itemize}}}

\vspace{3pt}

% --- Made in Limburg ---
\item{\cventry{2011--2015}{Projectmanager Made in Limburg}{Mediahuis Belgi\"e}{Hasselt}{}{\vspace{1pt}
\begin{itemize}
    \item Bouwde Made in Limburg uit tot het grootste onafhankelijke B2B-medium van Limburg (110.000 unieke bezoekers/maand); ontwierp het revenue model en schaalde naar vijf bijkomende regio's.
\end{itemize}}}

\vspace{3pt}

% --- Eerdere ervaring ---
\item{\cventry{2004--2011}{Eerdere ervaring, sales en business development}{Vanerum Group \textperiodcentered\ Cegeka \textperiodcentered\ Concentra Media}{}{}{}}

\end{itemize}

% ============================================================
%     OPLEIDING, EXECUTIVE EDUCATION EN CERTIFICATEN
% ============================================================

\section{Opleiding, executive education en certificaten}
\vspace{1pt}
\begin{itemize}
\item Leading AI \& digital marketing strategy, INSEAD (2026)
\item Blue Ocean Strategy, INSEAD (2026)
\item Brand management, Vlerick Business School (2025)
\item Behavioural design, fundamentals en advanced, SUE Academy (2024--2025)
\item Claude Code in Action, een van zes Anthropic-certificaten (2026)
\item Think different with Duncan Wardle (2025) \textperiodcentered\ Verbaal Meesterschap, Remco Claassen (2024)
\item Advanced automotive management, Febiac Academy (2020--2021)
\item Bachelor corporate communication, XIOS Hogeschool Limburg, nu PXL (2001--2004)
\end{itemize}

% ============================================================
%     TALEN
% ============================================================

\section{Talen}
\vspace{1pt}
\begin{itemize}
\item Nederlands (moedertaal), Engels (professioneel), Frans (professioneel).
\end{itemize}

% ============================================================
%     REFERENTIES
% ============================================================

\section{Referenties}
\vspace{1pt}
\begin{itemize}
\item{Beschikbaar op aanvraag.}
\end{itemize}

\end{document}
